\begin{prob}
    % Write the problem here.
    You are starting a new job, but your employer doesn't have much money to 
    pay you right now. Instead, they offer you two different payment plans:

    \textsc{Plan A}: On the first and second days, you'll be paid \$1. Every day
    thereafter, your pay will be \textbf{triple} the amount that it was two days previous.
    So your pay on day three is \$3, your pay on day 4 is \$3, your pay on day 5 is \$9, and so on.

    \textsc{Plan B}: On the first day, you'll be paid \$1, and your pay will double
    every day thereafter. On the second day you'll be paid \$2, on the third you'll
    be paid \$4, and so on.

    Which payment plan should you select if your goal is to make as much money as
    possible, asymptotically?

    \begin{soln}
        % This is the solution.
        For plan $A$, the recurrence relation is:
        \[
            S_A(n) = \begin{cases}
                3 \cdot S_A(n-2),& n \geq 3\\
                1,& n =1,2
            \end{cases}
        \]

        Unrolling the recurrence several times, we find:
        \begin{align*}
            S_A(n) 
        &= 
        3 \cdot S_A(n-2)
        \\
        &= 
        9 \cdot S_A(n-4)
        \\
        &= 
        27 \cdot S_A(n-6)
        \end{align*}
        In general, it looks like $S_A(n) = 3^k \cdot S_A(n - 2k)$. 

        On what step does the unrolling reach the base case? There are two bases cases,
        $S(1)$ and $S(2)$. Which base case is reached depends on whether $n$ is even or odd
        (since the argument to $S_A$ is being decreased by two at every step). It won't
        matter which base case we use in the end, but we can be careful and solve for both.

        If $n$ is even, then the $S_A(2)$ base case is reached. Solving for the number
        of steps needed:
        \[
            n - 2k = 2 \implies k = \frac{n-2}{2}
        \]
        If $n$ is odd, then the $S_A(1)$ base case is reached. Solving for the number
        of steps needed:
        \[
            n - 2k = 1 \implies k = \frac{n-1}{2}
        \]

        Plugging this value into the general solution, we find that:
        \[
            S_A(n) = 
            \begin{cases}
                3^{(n+2)/2},& n \text{ is even}\\
                3^{(n+1)/2},& n \text{ is odd}
            \end{cases}
        \]
        Asymptotically, $S_A(n) = \Theta(3^{n/2})$, since, for instance,
        $3^{(n+2)/2} = 3^{n/2} \cdot 3^{1/2} = \Theta(3^{n/2})$.

        We saw in lecture that the second payment plan leads to a payment of $\Theta(2^n)$
        on day $n$.

        Now we compare $\Theta(3^{n/2})$ to $\Theta(2^n)$. Note that $3^{n/2} =
        \left(3^{1/2}\right)^n$, and $3^{1/2} < 2$. Therefore $2^n$ is
        significantly larger than $3^{n/2}$ when $n$ is large, and we should
        prefer payment plan B.
    \end{soln}
\end{prob}
